Integrasi Pembangkit Listrik Tenaga Bayu (PLTB) berbasis \textit{doubly fed
induction generator} (DFIG-WT) dan \textit{permanent magnet synchronous
generator} (PMSG-WT) mengubah karakteristik stabilitas sinyal kecil sistem
tenaga listrik dengan cara yang bergantung pada arsitektur konverter
masing-masing teknologi. Penelitian ini menganalisis dan membandingkan dampak
integrasi DFIG-WT dan PMSG-WT terhadap \textit{small signal stability} (SSS)
pada sistem uji IEEE \textit{nine-bus} melalui analisis modal berbasis
\textit{eigenvalue} di DIgSILENT PowerFactory 2024, mencakup sebelas skenario
yang memvariasikan teknologi PLTB, posisi integrasi (Bus 7 dan Bus 9), serta
kondisi pembebanan. Integrasi DFIG-WT menurunkan \textit{damping ratio} mode
dominan dari 7,22\% menjadi 4,04\% terlepas dari posisi integrasi, dan
mempersempit margin kestabilan beban: batas ketidakstabilan Beban A turun dari
247\% menjadi 83\%, Beban B dari 373\% menjadi 123\%, keduanya menghasilkan
ketidakstabilan aperiodik (DR $-$100\%) yang berbeda dari ketidakstabilan
osilatoris pada sistem konvensional. Integrasi PMSG-WT mempertahankan atau
meningkatkan \textit{damping ratio} ke 7,25\%--8,73\%, dengan batas Beban A
165\% dan Beban B 284\%, karena \textit{fully-rated converter} memisahkan
dinamika mesin dari jaringan sehingga mode osilasi dominan tetap dikendalikan
oleh generator sinkron konvensional yang tersisa. Analisis
\textit{participation factor} mengonfirmasi bahwa variabel poros DFIG-WT
($\Delta\delta_{12}$) mendominasi mode osilasi dominan, sementara sudut rotor
internal DFIG-WT ($\varphi_m$) mendominasi mode tidak stabil saat beban
meningkat. DFIG-WT memperburuk redaman dan mempersempit margin operasi aman;
PMSG-WT mempertahankan atau meningkatkannya, dengan besaran peningkatan
bergantung pada generator yang digantikan.

\noindent{Kata kunci} : stabilitas sinyal kecil, DFIG-WT, PMSG-WT, analisis modal, \textit{eigenvalue}
